% Group report, IEEE conference format. The divide-and-conquer half lives in
% sec_dc.tex / app_dc.tex; everything the greedy half still needs is marked
% with \partnertodo{...} and prints in red.
%
% To build on Overleaf: upload this folder, keep the file layout, compile with
% pdfLaTeX. Nothing needs installing.
\documentclass[conference]{IEEEtran}
\IEEEoverridecommandlockouts
\usepackage{cite}
\usepackage{amsmath,amssymb,amsfonts}
\usepackage{amsthm}
\usepackage{algorithmic}
\usepackage{graphicx}
\usepackage{textcomp}
\usepackage{xcolor}
\usepackage{url}
\usepackage{listings}

% Theorem-like environments with air around them. IEEEtran's defaults have
% no vertical space, which makes a run of lemmas and proofs hard to read.
\newtheoremstyle{spaced}
  {10pt plus 2pt minus 2pt}  % space above
  {10pt plus 2pt minus 2pt}  % space below
  {\itshape}                 % body font
  {}                         % indent
  {\itshape}                 % head font
  {:}                        % punctuation after head
  {.5em}                     % space after head
  {}                         % head spec (default)
\theoremstyle{spaced}
\newtheorem{lemma}{Lemma}
\newtheorem{theorem}{Theorem}
\newtheorem{corollary}{Corollary}
\renewenvironment{IEEEproof}[1][\IEEEproofname]{\par\vspace{4pt}%
  \noindent\hspace{\parindent}{\itshape #1: }\ignorespaces}%
  {\hfill\IEEEQED\par\vspace{6pt}}

% Code listings (Appendix A) and the Lean proof (Appendix B).
\lstset{basicstyle=\ttfamily\scriptsize, breaklines=true, columns=fullflexible,
  frame=single, keepspaces=true, showstringspaces=false,
  % Unicode symbols in the Lean listing, mapped to math symbols.
  literate={≤}{{$\le$}}1 {≠}{{$\ne$}}1 {∧}{{$\wedge$}}1 {∨}{{$\vee$}}1
           {¬}{{$\neg$}}1 {↔}{{$\leftrightarrow$}}1 {→}{{$\rightarrow$}}1
           {∀}{{$\forall$}}1 {∃}{{$\exists$}}1 {⟨}{{$\langle$}}1 {⟩}{{$\rangle$}}1
           {×}{{$\times$}}1 {α}{{$\alpha$}}1 {∈}{{$\in$}}1 {·}{{$\cdot$}}1
           {₁}{{$_1$}}1 {₂}{{$_2$}}1 {⊢}{{$\vdash$}}1
           {⌈}{{$\lceil$}}1 {⌉}{{$\rceil$}}1 {⌊}{{$\lfloor$}}1 {⌋}{{$\rfloor$}}1}

% Work still owed by the greedy half, printed in red so nothing is missed.
\newcommand{\partnertodo}[1]{\textcolor{red}{\textbf{[Partner: #1]}}}

% Folder holding the divide-and-conquer figures and code (empty here).
\newcommand{\dcdir}{}

\begin{document}

\title{Greedy and Divide-and-Conquer Algorithms for Real-World Problems}

\author{\IEEEauthorblockN{Ilan Danial}
\IEEEauthorblockA{\textit{Dept.\ of Computer \& Information Science \& Engineering} \\
\textit{University of Florida}\\
Gainesville, FL, USA \\
ilan.danial@ufl.edu}
\and
\IEEEauthorblockN{\partnertodo{name}}
\IEEEauthorblockA{\textit{Dept.\ of Computer \& Information Science \& Engineering} \\
\textit{University of Florida}\\
Gainesville, FL, USA \\
\partnertodo{email}}
}

\maketitle

\begin{abstract}
\partnertodo{one or two sentences on the greedy problem and result.}
Map-search sites cap the number of listings shown per window, so collecting
every listing in a city means choosing windows that never overflow. We model
this as a point set on an integer grid and give a divide-and-conquer algorithm
that recursively splits overflowing windows into quadrants. We prove it
correct and bound its recursion depth by $\lceil\log_2 W\rceil$ for a map of
$W$ grid units; every statement is machine-checked in Lean~4. For evenly
spread listings it makes $\Theta(n/k)$ requests and $\Theta(n\log(n/k))$
work, which experiments confirm on synthetic data and on 261 real Airbnb
listings.
\end{abstract}

\begin{IEEEkeywords}
\partnertodo{greedy keywords}, divide and conquer, quadtree, web scraping,
formal verification
\end{IEEEkeywords}

\section{Introduction}

This report applies two algorithm-design techniques to problems that arose
outside the classroom. For each we state the problem, abstract it, give an
algorithm with a running-time analysis and a proof of correctness, explain the
algorithm in the problem's own terms, and check the analysis experimentally.

\partnertodo{one paragraph introducing the greedy problem and summarizing
the result.}

The second problem comes from web data collection. A map-search service
shows at most $k$ listings for any map window, so a scraper that needs every
listing in a city must split the map into windows that never overflow, using
as few requests as it can. Section~\ref{sec:dc} models the map as an integer
grid and the service as a capped oracle, solves the problem with a
divide-and-conquer recursion on quadrants, proves the recursion correct and
bounds its depth by $\lceil\log_2 W\rceil$, and verifies the predicted
$\Theta(n\log(n/k))$ cost experimentally. Every statement in the correctness
section is also machine-checked in Lean~4 (Appendix~\ref{app:dc-lean}).

\section{Greedy: \partnertodo{section title}}
\label{sec:greedy}

\partnertodo{the greedy half: problem, abstraction, algorithm, running time,
proof of correctness, explanation in the problem domain, experiments. Code
goes in an appendix after Appendix~\ref{app:dc-lean}; LLM prompts and results
in an appendix after Appendix~\ref{app:dc-llm}.}

% Divide-and-conquer half of the report.
% Preamble needs: amsmath, amssymb, algorithmic, graphicx, cite, and
%   \newtheorem{lemma}{Lemma}  \newtheorem{theorem}{Theorem}
% \dcdir is the folder holding figures/ (empty when compiled from this folder).

\section{Divide and Conquer: Capped Map Search}
\label{sec:dc}

\subsection{Problem}
\label{sec:dc-problem}

Airbnb's map search returns the listings inside the visible map window, but
never more than a fixed number of them: a crowded window is silently
truncated. A market study of a city needs \emph{every} listing, so a scraper
must choose map windows small enough that none is truncated, using as few
searches as possible. This problem came up while scraping Miami, FL and
Rinc\'on, PR~\cite{danial2023airbnb}; a single city-wide window returned only a
small fraction of the listings.

\subsection{Abstraction}
\label{sec:dc-abstraction}

Let $P$ be a set of $n$ \emph{listings}, each with a location in
$\mathbb{R}^2$. Several listings may share a location, since Airbnb rounds
coordinates; let $m$ be the largest number sharing one. A \emph{window} is a
half-open rectangle $r=[x_0,x_1)\times[y_0,y_1)$, and $P_r$ denotes the
listings located in $r$. The search region $R$ is a window with $P_R=P$.

The search service is an oracle with cap $k\ge1$:
\begin{equation}
Q_k(r)=\begin{cases}
P_r & \text{if } |P_r|\le k,\\
\textsc{overflow} & \text{otherwise.}
\end{cases}
\label{eq:oracle}
\end{equation}
The problem is to compute $P$ from $R$ and $Q_k$ using as few oracle calls
(web requests) as possible.

Splitting a window at its midpoint, both horizontally and vertically, gives
four equal \emph{quadrants}. Splitting recursively yields a
\emph{quadtree}~\cite{finkel1974quad,samet1984quadtree} whose leaves are the
windows actually read.

\subsection{Algorithm}
\label{sec:dc-algorithm}

\textsc{Collect} (Fig.~\ref{fig:dc-alg}) queries a window. If it is not
truncated, its listings are the answer. Otherwise the window is divided into
four quadrants, each is solved recursively, and the answers are united. The
full answer is $\textsc{Collect}(R)$.

\begin{figure}[t]
\begin{algorithmic}[1]
\STATE \textbf{function} \textsc{Collect}$(r)$
\STATE $A\gets Q_k(r)$
\IF{$A\neq\textsc{overflow}$}
  \RETURN $A$ \COMMENT{window fits under the cap}
\ENDIF
\STATE $r_1,\dots,r_4\gets$ quadrants of $r$ \COMMENT{divide}
\RETURN $\bigcup_{i=1}^{4}\textsc{Collect}(r_i)$ \COMMENT{conquer, combine}
\end{algorithmic}
\caption{Enumerating the listings in window $r$ through a capped oracle $Q_k$.}
\label{fig:dc-alg}
\end{figure}

% Placed here in the source so LaTeX can set it at the top of the next page;
% a figure* is always deferred to the page after the one where it appears.
\begin{figure*}[!t]
\centering
\includegraphics[width=\textwidth]{\dcdir figures/scaling}
\caption{Scaling with $n$ at $k=50$, for uniform and clustered listings.
(a)~Running time with the least-squares fit $c\,n\log_2(n/k)$
($R^2=0.995$ uniform, $0.996$ clustered). (b)~Running time divided by
$n\log_2(n/k)$; a flat curve confirms the growth rate. (c)~Requests issued,
against the lower bound $n/k$.}
\label{fig:dc-scaling}
\end{figure*}

\subsection{Correctness}
\label{sec:dc-correctness}

\begin{lemma}\label{lem:partition}
The four quadrants of a window $r$ partition $r$.
\end{lemma}
\begin{IEEEproof}
Let $(x_m,y_m)$ be the midpoint of $r$. For $(x,y)\in r$, exactly one of
$x<x_m$ and $x\ge x_m$ holds, and exactly one of $y<y_m$ and $y\ge y_m$; the
pair picks exactly one quadrant. Each quadrant lies in $r$ since
$x_0\le x_m\le x_1$ and $y_0\le y_m\le y_1$.
\end{IEEEproof}

Half-open windows are what make this work: a listing on a split line falls in
exactly one quadrant, so it is neither lost nor counted twice.

Scale $R$ to the unit square and let $\delta\in(0,1]$ be the smallest
coordinate-wise distance, $\max(|\Delta x|,|\Delta y|)$, between two listings at
different locations ($\delta=1$ if there is only one location).

\begin{lemma}\label{lem:height}
If $m\le k$, the recursion tree of $\textsc{Collect}(R)$ has height at most
$\lceil\log_2(1/\delta)\rceil$; in particular the algorithm terminates.
\end{lemma}
\begin{IEEEproof}
A window at depth $d$ has side $2^{-d}$, so two listings inside it at different
locations are less than $2^{-d}$ apart, i.e.\ $\delta<2^{-d}$. A window is split
only if it holds more than $k\ge m$ listings, hence listings at two or more
locations, hence $d<\log_2(1/\delta)$. So every split window has depth at most
$\lceil\log_2(1/\delta)\rceil-1$ and every window depth at most
$\lceil\log_2(1/\delta)\rceil$. A tree of bounded height with four children per
internal node is finite.
\end{IEEEproof}

\begin{theorem}\label{thm:correct}
If $m\le k$, $\textsc{Collect}(r)=P_r$ for every window $r$, and the leaf
windows partition $r$ with at most $k$ listings each.
\end{theorem}
\begin{IEEEproof}
Strong induction on the height $h$ of the recursion tree at $r$ (finite by
Lemma~\ref{lem:height}). If $h=0$, the oracle did not overflow, so
$|P_r|\le k$ and line~4 returns $Q_k(r)=P_r$; the only leaf is $r$. If $h>0$,
the algorithm recurses on the quadrants $r_1,\dots,r_4$, whose subtrees are
shorter than $h$, so by induction $\textsc{Collect}(r_i)=P_{r_i}$ and the leaves
under $r_i$ partition $r_i$ with at most $k$ listings each. By
Lemma~\ref{lem:partition} each listing of $P_r$ lies in exactly one $r_i$, so
$\bigcup_i P_{r_i}=P_r$ and the leaves of the four subtrees partition $r$.
\end{IEEEproof}

The condition $m\le k$ is also necessary: if more than $k$ listings share a
location, every window containing it overflows, so no strategy using $Q_k$
can reveal them.

Lemmas~\ref{lem:partition}--\ref{lem:height} and Theorem~\ref{thm:correct}
are also machine-checked in Lean~4~\cite{lean4} (Appendix~\ref{app:dc-lean}).
Reported coordinates have finite precision (Airbnb gives five decimals), so
the formalization works on an integer grid, where $\delta$ is one grid unit.
There Lemma~\ref{lem:height} reads: if $m\le k$, recursion depth
$\lceil\log_2 W\rceil$ suffices for a window of side $W$; combined with
Theorem~\ref{thm:correct} this gives total correctness. The real-valued form
of Lemma~\ref{lem:height} is proved only on paper. The nearest existing
formalizations, region quadtrees in Isabelle~\cite{nipkow2024quad,nipkow2024afp}
and $k$-d trees in Isabelle and Coq~\cite{rau2019kdtree,hamid2023kdtree},
obtain termination from a fixed grid exponent or from the point count
shrinking at each median split; neither covers a space-splitting tree over
arbitrary points, so the proof here was written from scratch.

\subsection{Running Time}
\label{sec:dc-analysis}

Let $N$ be the number of oracle calls and $T$ the total work when a node with
$|P_r|$ listings spends $\Theta(|P_r|+1)$ time returning or distributing them.
Let $h$ be the tree height.

Windows at one depth are disjoint (Lemma~\ref{lem:partition}), so each level
holds at most $n$ listings and at most $\lfloor n/(k+1)\rfloor$ split windows,
each holding at least $k+1$ listings. Internal nodes sit at depths
$0,\dots,h-1$ and have four children each, so
\begin{equation}
\lceil n/k\rceil\le N\le 1+4h\Big\lfloor\frac{n}{k+1}\Big\rfloor,
\qquad T=O\big(n(h+1)\big).
\label{eq:general}
\end{equation}
The lower bound holds because the leaves cover $P$ with at most $k$ listings
each. By Lemma~\ref{lem:height}, $T=O(n\log(1/\delta))$: the cost is governed
by how tightly listings crowd, not by $n$ alone.

In the balanced case, where each split sends about $n/4$ listings to each
quadrant (true in expectation for uniformly random locations),
\begin{equation}
T(n)=4\,T(n/4)+\Theta(n)\quad(n>k),
\label{eq:recurrence}
\end{equation}
with $T(n)=\Theta(1)$ for $n\le k$. With $a=b=4$ and
$f(n)=\Theta(n)=\Theta(n^{\log_b a})$, case~2 of the Master
Theorem~\cite{cormen2022introduction} gives $T(n)=\Theta(n\log n)$. Unrolling
the recursion shows the dependence on $k$: the tree is complete with
$L=\lceil\log_4(n/k)\rceil$ levels below the root, each costing $\Theta(n)$, so
\begin{align}
T(n)&=\Theta\!\left(n\log\frac{n}{k}\right),\label{eq:balanced}\\
N&=\sum_{i=0}^{L}4^i=\frac{4^{L+1}-1}{3}=\Theta\!\left(\frac{n}{k}\right),
\label{eq:requests}
\end{align}
within a constant factor of the $n/k$ requests any strategy must make.

\subsection{In the Language of the Problem}
\label{sec:dc-domain}

Search the whole city in one window. If the results fit on the pages the site
shows, record them. Otherwise cut the window into four equal tiles and repeat
on each. A listing exactly on a cut goes to the tile right of or above it, so
it is recorded once. Dense neighborhoods get zoomed into repeatedly; open water
is dismissed after one search. The number of searches is proportional to
$n/k$, the number of listings divided by the number shown per search, by
\eqref{eq:requests}. The only obstacle is an address with more listings than
the cap, such as an apartment building whose units share rounded coordinates:
no zoom level separates them, which is why Theorem~\ref{thm:correct} needs
$m\le k$. In Rinc\'on at most $7$ listings share a location, so any $k\ge7$
works.

\subsection{Experiments}
\label{sec:dc-experiments}

\begin{figure}[!t]
\centering
\includegraphics[width=\columnwidth]{\dcdir figures/depth_vs_delta}
\caption{Effect of packing $k+1$ listings into a square of side $\delta$
($n=10{,}000$, $k=50$): (a)~tree height, (b)~requests.}
\label{fig:dc-depth}
\end{figure}

\begin{figure}[!t]
\centering
\includegraphics[width=2.45in]{\dcdir figures/rincon_quadtree}
\caption{Windows chosen for 261 listings in Rinc\'on, PR ($k=20$).}
\label{fig:dc-rincon}
\end{figure}

\textbf{Setup.} \textsc{Collect} was implemented in Python~3.14
(Appendix~\ref{app:dc-code}) and run on an AMD Ryzen~7 9800X3D. The oracle is
simulated: each node partitions the listings it holds among its quadrants,
which is the $\Theta(|P_r|)$ step of \eqref{eq:recurrence}. We used $k=50$
and $n=2^{10},\dots,2^{19}$ listings in the unit square, either uniform or
drawn from $20$ Gaussian clusters ($\sigma=0.02$). Each size was run five
times on fresh random inputs; we report the fastest run with garbage
collection disabled, as \texttt{timeit} does~\cite{python-timeit}. Every run
was checked to return each listing exactly once.

\textbf{Running time.} Fig.~\ref{fig:dc-scaling}(a) shows measured time
against the fitted model; Fig.~\ref{fig:dc-scaling}(b) divides time by
$n\log_2(n/k)$. Over a $512\times$ range of $n$ the ratio stays within
$26$--$43$\,ns (uniform) and $36$--$48$\,ns (clustered), confirming
$\Theta(n\log(n/k))$. The zig-zag arises because the tree gains a level only
when $n$ grows about fourfold, so the true cost rises in steps; the slight rise
at the largest sizes is consistent with cache effects. Clustered inputs cost
about $20\%$ more because their trees are two levels taller.

\textbf{Requests.} Table~\ref{tab:dc-requests} compares requests on uniform
inputs with the exact balanced count \eqref{eq:requests}; they agree within
$3.1\%$. Fig.~\ref{fig:dc-scaling}(c) shows the count growing linearly,
parallel to the lower bound $n/k$.

\begin{table}[t]
\caption{Requests on uniform inputs ($k=50$) vs.\ the prediction \eqref{eq:requests}}
\label{tab:dc-requests}
\centering
\begin{tabular}{rrrr}
\hline
$n$ & measured & predicted & height \\
\hline
1{,}024   & 85     & 85     & 3 \\
4{,}096   & 333    & 341    & 4 \\
16{,}384  & 1{,}337  & 1{,}365  & 5 \\
65{,}536  & 5{,}293  & 5{,}461  & 6 \\
262{,}144 & 21{,}197 & 21{,}845 & 7 \\
524{,}288 & 21{,}917 & 21{,}845 & 8 \\
\hline
\end{tabular}
\end{table}

\textbf{Crowding.} To test \eqref{eq:general}, we fixed $n=10{,}000$ uniform
listings and added $k+1=51$ listings inside a square of side
$\delta=2^{-2},\dots,2^{-28}$. Fig.~\ref{fig:dc-depth} shows the tree height
tracking $\log_2(1/\delta)$ (Lemma~\ref{lem:height}) although $n$ is fixed,
and each extra level adding about four requests, the $4h$ term of
\eqref{eq:general}.

\textbf{Real data.} Fig.~\ref{fig:dc-rincon} runs the algorithm with $k=20$ on
$261$ Airbnb listings from Rinc\'on, PR (scraped March 29, 2023~\cite{danial2023airbnb}).
It used $57$ requests, $43$ windows and height $5$, within the bounds
$14\le N\le 241$ of \eqref{eq:general}. The two dense clusters are split
repeatedly; the outskirts are covered by a few large windows.


\section{Conclusion}

\partnertodo{one paragraph concluding the greedy half.}

For the capped map search, splitting overflowing windows into quadrants
collects every listing with a number of requests within a constant factor of
the $n/k$ lower bound when listings are spread evenly, and never recurses
deeper than $\lceil\log_2 W\rceil$ levels on a map of $W$ grid units. The
one assumption, that no single location holds more than $k$ listings, is
necessary for any strategy that only sees capped search results. Measured
running times match the $\Theta(n\log(n/k))$ analysis within a few percent
over three orders of magnitude, and the correctness argument is
machine-checked.

\bibliographystyle{IEEEtran}
\bibliography{refs}

% Code reads better at full page width.
\onecolumn
\appendices
% Appendices for the divide-and-conquer half. Needs \usepackage{listings}.

\section{Divide-and-Conquer Code}
\label{app:dc-code}

The algorithm (\texttt{quadtree.py}) and the measurements behind
Figs.~\ref{fig:dc-scaling}--\ref{fig:dc-depth} and Table~\ref{tab:dc-requests}
(\texttt{experiments.py}). The plotting script, the unit tests, the Rinc\'on
data file and the raw measurements are in the project
repository.\footnote{\url{https://github.com/IlanDanial/P1Alg}}

\lstinputlisting[language=Python,caption={quadtree.py}]{\dcdir quadtree.py}
\lstinputlisting[language=Python,caption={experiments.py}]{\dcdir experiments.py}

\section{Machine-Checked Proof}
\label{app:dc-lean}

Lemma~\ref{lem:partition} (\texttt{quadrants\_cover},
\texttt{quadrants\_disjoint}), Theorem~\ref{thm:correct}
(\texttt{collect\_correct}), Lemma~\ref{lem:height} on the integer grid
(\texttt{collect\_terminates}) and their combination into total correctness
(\texttt{collect\_total}) in Lean~4.35, using only the core library.
\texttt{lake build} in \texttt{divide\_and\_conquer/proof} checks it; the
theorems depend on no axioms beyond Lean's standard three
(\texttt{propext}, \texttt{Classical.choice}, \texttt{Quot.sound}).

\lstinputlisting[language={},caption={proof/Proof/Collect.lean}]{\dcdir proof/Proof/Collect.lean}

\section{LLM Usage: Divide-and-Conquer Part}
\label{app:dc-llm}

\textbf{Tool.} Claude Code (Anthropic's command-line coding agent), model
Claude Opus~5.5 (\texttt{claude-opus-5-5}), October 2026. The agent read the
assignment, read the scraper repository~\cite{danial2023airbnb}, searched the web
for the IEEE template, wrote and ran the code, and drafted this section. All
proofs, numbers, and citations were checked by the authors.

\textbf{Prompts (verbatim).}
\begin{enumerate}
\item ``DO you understand the project I need to do also this is a group project
  I am doing the divide an conquer part''
\item ``My Idea for the problem is a problem I have solved before with divide
  and conquer \url{https://github.com/IlanDanial/airbnb_data} I split the map over and
  over, each time checking the amount of properties in the split and if it was
  under I took those properties and added it to the list since you can only see
  so many properties at a time''
\item ``yes start the implementation, we're using IEEE search online for proper
  template''
\item ``yes write the LaTeX section''
\item ``This is really ugly'' followed by ``My bigger issue is that the
  explanations are very longwinded for no reason and the formatting of the
  page is ugly its hard to read as well''
\item ``I have an issue with some of this read the paper be critical how do
  you think of the formatting'', then ``fix all of them at once'' and ``also
  do we need the full code in the appendix or can we reference a repository''
\item ``would it be worth it to look at the code and refactor it and the
  comments in the code'', then ``go''
\item ``Isnt there a code languge that is built to prove proofs can we use
  that to prove correctness'', then ``I want to go for it also I am sure there
  are online proofs of this exact same thing in the same language so it should
  take less than an hour''
\item ``so your telling me that there was no problem with proving the whole
  thing end to end''
\item ``actually search online for the hard part dont just tell me it does
  not exist''
\item ``do both, go for the full proof, or is there a solution we can do that
  would be easier to prove and solve the problem all the same''
\end{enumerate}

\textbf{Results.}
% TODO: paste the full exported transcript (Claude Code: /export) here or in
% supplementary material; the assignment requires all intermediate results.
(1)~A summary of the assignment and four candidate divide-and-conquer problems.
(2)~An analysis of the scraper's \texttt{subdivide} recursion as an adaptive
quadtree, identifying the termination issue for co-located listings and the
double-counting issue on window borders.
(3)~\texttt{quadtree.py}, \texttt{test\_quadtree.py}, \texttt{experiments.py},
and the figures; an intermediate run stopped when the minimum-width guard was
reached, after which the crowding experiment was limited to
$\delta\ge2^{-28}$, and garbage collection was disabled during timing.
(4)~A first draft of Section~\ref{sec:dc} and these appendices, with five
separate figures in matplotlib's default style.
(5)~The draft was shortened by about a third, the five figures were
restyled to match the paper's typeface and merged into three, and the page
layout was fixed.
(6)~A list of nine formatting defects (an equation label pushed to its own
line, a wrapped pseudocode line, mixed paragraph styles, inconsistent
equation citations, a legend mismatch), all then fixed; the unit tests were
removed from this appendix in favor of the repository link.
(7)~The plotting code was split out of \texttt{experiments.py} into
\texttt{plots.py}, and docstrings and comments were added.
(8)~An assessment of proof assistants (Lean~4, Coq, Isabelle, Dafny) and of
which parts of the proof are feasible to formalize; then a first Lean
formalization of Lemma~\ref{lem:partition} and Theorem~\ref{thm:correct}
only, with \textsc{Collect} given a fuel bound, which compiled on the first
build.
(9)~A clarification that this was partial correctness, not an end-to-end
proof, with three further gaps listed (integer midpoints, the oracle not being
a black box, no duplicate-freeness proof).
(10)~A web search that found the Isabelle and Coq formalizations cited in
Section~\ref{sec:dc-correctness}, none covering separation-based termination;
the assistant's earlier claim that no related formalization existed was
withdrawn.
(11)~The observation that finite-precision coordinates make the problem a
grid problem, on which Lemma~\ref{lem:height} is provable by induction on the
depth; \texttt{collect\_terminates} and \texttt{collect\_total} were then
added, needing one round of fixes to core-library lemma names before
compiling.


\end{document}
